% Shared Beamer preamble for Introduction to Programming with Kotlin
% Include with: % Shared Beamer preamble for Introduction to Programming with Kotlin
% Include with: % Shared Beamer preamble for Introduction to Programming with Kotlin
% Include with: % Shared Beamer preamble for Introduction to Programming with Kotlin
% Include with: \input{preamble}

\usetheme{Madrid}
\usecolortheme{default}

% Clean, professional palette (deep navy + muted teal accent)
\definecolor{LectureNavy}{HTML}{1B3A4B}
\definecolor{LectureAccent}{HTML}{2A9D8F}
\definecolor{LectureGray}{HTML}{4A5568}
\definecolor{CodeBg}{HTML}{F7FAFC}
\definecolor{AlertBg}{HTML}{FFF5F5}
\definecolor{TryBg}{HTML}{F0FFF4}
\definecolor{QuizBg}{HTML}{EBF8FF}
\definecolor{AnswerKeyBg}{HTML}{FFFFF0}
\definecolor{AnswerGold}{HTML}{B7791F}
\definecolor{KeywordBlue}{HTML}{2B6CB0}
\definecolor{StringGreen}{HTML}{276749}
\definecolor{CommentGray}{HTML}{718096}

\setbeamercolor{structure}{fg=LectureNavy}
\setbeamercolor{palette primary}{bg=LectureNavy,fg=white}
\setbeamercolor{palette secondary}{bg=LectureAccent,fg=white}
\setbeamercolor{palette tertiary}{bg=LectureGray,fg=white}
% Madrid puts title/subtitle in a coloured box (inherits palette primary bg).
% Title text must be light on that dark box — not navy-on-navy.
\setbeamercolor{title}{fg=white,bg=LectureNavy}
\setbeamercolor{subtitle}{fg=white,bg=LectureNavy}
\setbeamercolor{author}{fg=LectureNavy}
\setbeamercolor{date}{fg=LectureGray}
\setbeamercolor{institute}{fg=LectureGray}
\setbeamercolor{frametitle}{bg=LectureNavy,fg=white}
\setbeamercolor{block title}{bg=LectureNavy,fg=white}
\setbeamercolor{block body}{bg=CodeBg,fg=black}
\setbeamercolor{block title alerted}{bg=LectureAccent,fg=white}
\setbeamercolor{block body alerted}{bg=TryBg,fg=black}
\setbeamercolor{block title example}{bg=LectureGray,fg=white}
\setbeamercolor{block body example}{bg=AlertBg,fg=black}

\setbeamertemplate{navigation symbols}{}
\setbeamertemplate{itemize items}[circle]
\setbeamertemplate{enumerate items}[default]

\usepackage[T1]{fontenc}
\usepackage[utf8]{inputenc}
\usepackage{lmodern}
\usepackage{booktabs}
\usepackage{tabularx}
\usepackage{graphicx}
\usepackage{hyperref}
\usepackage{listings}
\usepackage{xcolor}
\usepackage{tikz}
\usetikzlibrary{positioning,arrows.meta}

% listings: Kotlin without shell-escape (no built-in Kotlin language)
\lstdefinelanguage{Kotlin}{
  morekeywords={
    abstract,actual,annotation,as,break,by,catch,class,companion,const,
    constructor,continue,crossinline,data,do,dynamic,else,enum,expect,
    external,false,field,file,final,finally,for,fun,get,if,import,in,
    infix,init,inline,inner,interface,internal,is,lateinit,noinline,null,
    object,open,operator,out,override,package,private,protected,public,
    reified,return,sealed,set,super,suspend,tailrec,this,throw,true,try,
    typealias,typeof,val,var,when,where,while
  },
  sensitive=true,
  morecomment=[l]{//},
  morecomment=[s]{/*}{*/},
  morestring=[b]",
  morestring=[b]',
  morestring=[s]{"""}{"""},
}

\lstdefinestyle{kotlinstyle}{
  language=Kotlin,
  basicstyle=\ttfamily\small,
  keywordstyle=\color{KeywordBlue}\bfseries,
  commentstyle=\color{CommentGray}\itshape,
  stringstyle=\color{StringGreen},
  numbers=left,
  numberstyle=\tiny\color{CommentGray},
  numbersep=6pt,
  frame=single,
  rulecolor=\color{LectureGray!40},
  backgroundcolor=\color{CodeBg},
  breaklines=true,
  showstringspaces=false,
  tabsize=2,
  captionpos=b,
  morekeywords={readln,readLine,println,print,listOf,mutableListOf,toInt,toDouble,toLong,toShort,toByte,toChar,toString,split,size},
}
\lstset{style=kotlinstyle}

\newcommand{\kotlincode}[1]{\texttt{\small #1}}
\newcommand{\trythis}[1]{%
  \begin{alertblock}{Try this}
    #1
  \end{alertblock}%
}
\newcommand{\commonmistake}[1]{%
  \begin{exampleblock}{Common mistake}
    #1
  \end{exampleblock}%
}
% Formative in-class checks: question slide, then a dedicated Answer slide
\newcommand{\quizpause}{%
  \par\vspace{0.5em}
  {\small\textit{Pause --- answer (clicker / raise hand / neighbour) before the next slide.}}%
}
\newcommand{\quizbox}[1]{%
  \setbeamercolor{block title}{bg=KeywordBlue,fg=white}%
  \setbeamercolor{block body}{bg=QuizBg,fg=black}%
  \begin{block}{Quiz}
    #1
  \end{block}%
  \setbeamercolor{block title}{bg=LectureNavy,fg=white}%
  \setbeamercolor{block body}{bg=CodeBg,fg=black}%
}
\newcommand{\answerbox}[1]{%
  \setbeamercolor{block title}{bg=AnswerGold,fg=white}%
  \setbeamercolor{block body}{bg=AnswerKeyBg,fg=black}%
  \begin{block}{Answer}
    #1
  \end{block}%
  \setbeamercolor{block title}{bg=LectureNavy,fg=white}%
  \setbeamercolor{block body}{bg=CodeBg,fg=black}%
}

% Empty author/date placeholders for the lecturer to fill
\author{}
\date{}


\usetheme{Madrid}
\usecolortheme{default}

% Clean, professional palette (deep navy + muted teal accent)
\definecolor{LectureNavy}{HTML}{1B3A4B}
\definecolor{LectureAccent}{HTML}{2A9D8F}
\definecolor{LectureGray}{HTML}{4A5568}
\definecolor{CodeBg}{HTML}{F7FAFC}
\definecolor{AlertBg}{HTML}{FFF5F5}
\definecolor{TryBg}{HTML}{F0FFF4}
\definecolor{QuizBg}{HTML}{EBF8FF}
\definecolor{AnswerKeyBg}{HTML}{FFFFF0}
\definecolor{AnswerGold}{HTML}{B7791F}
\definecolor{KeywordBlue}{HTML}{2B6CB0}
\definecolor{StringGreen}{HTML}{276749}
\definecolor{CommentGray}{HTML}{718096}

\setbeamercolor{structure}{fg=LectureNavy}
\setbeamercolor{palette primary}{bg=LectureNavy,fg=white}
\setbeamercolor{palette secondary}{bg=LectureAccent,fg=white}
\setbeamercolor{palette tertiary}{bg=LectureGray,fg=white}
% Madrid puts title/subtitle in a coloured box (inherits palette primary bg).
% Title text must be light on that dark box — not navy-on-navy.
\setbeamercolor{title}{fg=white,bg=LectureNavy}
\setbeamercolor{subtitle}{fg=white,bg=LectureNavy}
\setbeamercolor{author}{fg=LectureNavy}
\setbeamercolor{date}{fg=LectureGray}
\setbeamercolor{institute}{fg=LectureGray}
\setbeamercolor{frametitle}{bg=LectureNavy,fg=white}
\setbeamercolor{block title}{bg=LectureNavy,fg=white}
\setbeamercolor{block body}{bg=CodeBg,fg=black}
\setbeamercolor{block title alerted}{bg=LectureAccent,fg=white}
\setbeamercolor{block body alerted}{bg=TryBg,fg=black}
\setbeamercolor{block title example}{bg=LectureGray,fg=white}
\setbeamercolor{block body example}{bg=AlertBg,fg=black}

\setbeamertemplate{navigation symbols}{}
\setbeamertemplate{itemize items}[circle]
\setbeamertemplate{enumerate items}[default]

\usepackage[T1]{fontenc}
\usepackage[utf8]{inputenc}
\usepackage{lmodern}
\usepackage{booktabs}
\usepackage{tabularx}
\usepackage{graphicx}
\usepackage{hyperref}
\usepackage{listings}
\usepackage{xcolor}
\usepackage{tikz}
\usetikzlibrary{positioning,arrows.meta}

% listings: Kotlin without shell-escape (no built-in Kotlin language)
\lstdefinelanguage{Kotlin}{
  morekeywords={
    abstract,actual,annotation,as,break,by,catch,class,companion,const,
    constructor,continue,crossinline,data,do,dynamic,else,enum,expect,
    external,false,field,file,final,finally,for,fun,get,if,import,in,
    infix,init,inline,inner,interface,internal,is,lateinit,noinline,null,
    object,open,operator,out,override,package,private,protected,public,
    reified,return,sealed,set,super,suspend,tailrec,this,throw,true,try,
    typealias,typeof,val,var,when,where,while
  },
  sensitive=true,
  morecomment=[l]{//},
  morecomment=[s]{/*}{*/},
  morestring=[b]",
  morestring=[b]',
  morestring=[s]{"""}{"""},
}

\lstdefinestyle{kotlinstyle}{
  language=Kotlin,
  basicstyle=\ttfamily\small,
  keywordstyle=\color{KeywordBlue}\bfseries,
  commentstyle=\color{CommentGray}\itshape,
  stringstyle=\color{StringGreen},
  numbers=left,
  numberstyle=\tiny\color{CommentGray},
  numbersep=6pt,
  frame=single,
  rulecolor=\color{LectureGray!40},
  backgroundcolor=\color{CodeBg},
  breaklines=true,
  showstringspaces=false,
  tabsize=2,
  captionpos=b,
  morekeywords={readln,readLine,println,print,listOf,mutableListOf,toInt,toDouble,toLong,toShort,toByte,toChar,toString,split,size},
}
\lstset{style=kotlinstyle}

\newcommand{\kotlincode}[1]{\texttt{\small #1}}
\newcommand{\trythis}[1]{%
  \begin{alertblock}{Try this}
    #1
  \end{alertblock}%
}
\newcommand{\commonmistake}[1]{%
  \begin{exampleblock}{Common mistake}
    #1
  \end{exampleblock}%
}
% Formative in-class checks: question slide, then a dedicated Answer slide
\newcommand{\quizpause}{%
  \par\vspace{0.5em}
  {\small\textit{Pause --- answer (clicker / raise hand / neighbour) before the next slide.}}%
}
\newcommand{\quizbox}[1]{%
  \setbeamercolor{block title}{bg=KeywordBlue,fg=white}%
  \setbeamercolor{block body}{bg=QuizBg,fg=black}%
  \begin{block}{Quiz}
    #1
  \end{block}%
  \setbeamercolor{block title}{bg=LectureNavy,fg=white}%
  \setbeamercolor{block body}{bg=CodeBg,fg=black}%
}
\newcommand{\answerbox}[1]{%
  \setbeamercolor{block title}{bg=AnswerGold,fg=white}%
  \setbeamercolor{block body}{bg=AnswerKeyBg,fg=black}%
  \begin{block}{Answer}
    #1
  \end{block}%
  \setbeamercolor{block title}{bg=LectureNavy,fg=white}%
  \setbeamercolor{block body}{bg=CodeBg,fg=black}%
}

% Empty author/date placeholders for the lecturer to fill
\author{}
\date{}


\usetheme{Madrid}
\usecolortheme{default}

% Clean, professional palette (deep navy + muted teal accent)
\definecolor{LectureNavy}{HTML}{1B3A4B}
\definecolor{LectureAccent}{HTML}{2A9D8F}
\definecolor{LectureGray}{HTML}{4A5568}
\definecolor{CodeBg}{HTML}{F7FAFC}
\definecolor{AlertBg}{HTML}{FFF5F5}
\definecolor{TryBg}{HTML}{F0FFF4}
\definecolor{QuizBg}{HTML}{EBF8FF}
\definecolor{AnswerKeyBg}{HTML}{FFFFF0}
\definecolor{AnswerGold}{HTML}{B7791F}
\definecolor{KeywordBlue}{HTML}{2B6CB0}
\definecolor{StringGreen}{HTML}{276749}
\definecolor{CommentGray}{HTML}{718096}

\setbeamercolor{structure}{fg=LectureNavy}
\setbeamercolor{palette primary}{bg=LectureNavy,fg=white}
\setbeamercolor{palette secondary}{bg=LectureAccent,fg=white}
\setbeamercolor{palette tertiary}{bg=LectureGray,fg=white}
% Madrid puts title/subtitle in a coloured box (inherits palette primary bg).
% Title text must be light on that dark box — not navy-on-navy.
\setbeamercolor{title}{fg=white,bg=LectureNavy}
\setbeamercolor{subtitle}{fg=white,bg=LectureNavy}
\setbeamercolor{author}{fg=LectureNavy}
\setbeamercolor{date}{fg=LectureGray}
\setbeamercolor{institute}{fg=LectureGray}
\setbeamercolor{frametitle}{bg=LectureNavy,fg=white}
\setbeamercolor{block title}{bg=LectureNavy,fg=white}
\setbeamercolor{block body}{bg=CodeBg,fg=black}
\setbeamercolor{block title alerted}{bg=LectureAccent,fg=white}
\setbeamercolor{block body alerted}{bg=TryBg,fg=black}
\setbeamercolor{block title example}{bg=LectureGray,fg=white}
\setbeamercolor{block body example}{bg=AlertBg,fg=black}

\setbeamertemplate{navigation symbols}{}
\setbeamertemplate{itemize items}[circle]
\setbeamertemplate{enumerate items}[default]

\usepackage[T1]{fontenc}
\usepackage[utf8]{inputenc}
\usepackage{lmodern}
\usepackage{booktabs}
\usepackage{tabularx}
\usepackage{graphicx}
\usepackage{hyperref}
\usepackage{listings}
\usepackage{xcolor}
\usepackage{tikz}
\usetikzlibrary{positioning,arrows.meta}

% listings: Kotlin without shell-escape (no built-in Kotlin language)
\lstdefinelanguage{Kotlin}{
  morekeywords={
    abstract,actual,annotation,as,break,by,catch,class,companion,const,
    constructor,continue,crossinline,data,do,dynamic,else,enum,expect,
    external,false,field,file,final,finally,for,fun,get,if,import,in,
    infix,init,inline,inner,interface,internal,is,lateinit,noinline,null,
    object,open,operator,out,override,package,private,protected,public,
    reified,return,sealed,set,super,suspend,tailrec,this,throw,true,try,
    typealias,typeof,val,var,when,where,while
  },
  sensitive=true,
  morecomment=[l]{//},
  morecomment=[s]{/*}{*/},
  morestring=[b]",
  morestring=[b]',
  morestring=[s]{"""}{"""},
}

\lstdefinestyle{kotlinstyle}{
  language=Kotlin,
  basicstyle=\ttfamily\small,
  keywordstyle=\color{KeywordBlue}\bfseries,
  commentstyle=\color{CommentGray}\itshape,
  stringstyle=\color{StringGreen},
  numbers=left,
  numberstyle=\tiny\color{CommentGray},
  numbersep=6pt,
  frame=single,
  rulecolor=\color{LectureGray!40},
  backgroundcolor=\color{CodeBg},
  breaklines=true,
  showstringspaces=false,
  tabsize=2,
  captionpos=b,
  morekeywords={readln,readLine,println,print,listOf,mutableListOf,toInt,toDouble,toLong,toShort,toByte,toChar,toString,split,size},
}
\lstset{style=kotlinstyle}

\newcommand{\kotlincode}[1]{\texttt{\small #1}}
\newcommand{\trythis}[1]{%
  \begin{alertblock}{Try this}
    #1
  \end{alertblock}%
}
\newcommand{\commonmistake}[1]{%
  \begin{exampleblock}{Common mistake}
    #1
  \end{exampleblock}%
}
% Formative in-class checks: question slide, then a dedicated Answer slide
\newcommand{\quizpause}{%
  \par\vspace{0.5em}
  {\small\textit{Pause --- answer (clicker / raise hand / neighbour) before the next slide.}}%
}
\newcommand{\quizbox}[1]{%
  \setbeamercolor{block title}{bg=KeywordBlue,fg=white}%
  \setbeamercolor{block body}{bg=QuizBg,fg=black}%
  \begin{block}{Quiz}
    #1
  \end{block}%
  \setbeamercolor{block title}{bg=LectureNavy,fg=white}%
  \setbeamercolor{block body}{bg=CodeBg,fg=black}%
}
\newcommand{\answerbox}[1]{%
  \setbeamercolor{block title}{bg=AnswerGold,fg=white}%
  \setbeamercolor{block body}{bg=AnswerKeyBg,fg=black}%
  \begin{block}{Answer}
    #1
  \end{block}%
  \setbeamercolor{block title}{bg=LectureNavy,fg=white}%
  \setbeamercolor{block body}{bg=CodeBg,fg=black}%
}

% Empty author/date placeholders for the lecturer to fill
\author{}
\date{}


\usetheme{Madrid}
\usecolortheme{default}

% Clean, professional palette (deep navy + muted teal accent)
\definecolor{LectureNavy}{HTML}{1B3A4B}
\definecolor{LectureAccent}{HTML}{2A9D8F}
\definecolor{LectureGray}{HTML}{4A5568}
\definecolor{CodeBg}{HTML}{F7FAFC}
\definecolor{AlertBg}{HTML}{FFF5F5}
\definecolor{TryBg}{HTML}{F0FFF4}
\definecolor{QuizBg}{HTML}{EBF8FF}
\definecolor{AnswerKeyBg}{HTML}{FFFFF0}
\definecolor{AnswerGold}{HTML}{B7791F}
\definecolor{KeywordBlue}{HTML}{2B6CB0}
\definecolor{StringGreen}{HTML}{276749}
\definecolor{CommentGray}{HTML}{718096}

\setbeamercolor{structure}{fg=LectureNavy}
\setbeamercolor{palette primary}{bg=LectureNavy,fg=white}
\setbeamercolor{palette secondary}{bg=LectureAccent,fg=white}
\setbeamercolor{palette tertiary}{bg=LectureGray,fg=white}
% Madrid puts title/subtitle in a coloured box (inherits palette primary bg).
% Title text must be light on that dark box — not navy-on-navy.
\setbeamercolor{title}{fg=white,bg=LectureNavy}
\setbeamercolor{subtitle}{fg=white,bg=LectureNavy}
\setbeamercolor{author}{fg=LectureNavy}
\setbeamercolor{date}{fg=LectureGray}
\setbeamercolor{institute}{fg=LectureGray}
\setbeamercolor{frametitle}{bg=LectureNavy,fg=white}
\setbeamercolor{block title}{bg=LectureNavy,fg=white}
\setbeamercolor{block body}{bg=CodeBg,fg=black}
\setbeamercolor{block title alerted}{bg=LectureAccent,fg=white}
\setbeamercolor{block body alerted}{bg=TryBg,fg=black}
\setbeamercolor{block title example}{bg=LectureGray,fg=white}
\setbeamercolor{block body example}{bg=AlertBg,fg=black}

\setbeamertemplate{navigation symbols}{}
\setbeamertemplate{itemize items}[circle]
\setbeamertemplate{enumerate items}[default]

\usepackage[T1]{fontenc}
\usepackage[utf8]{inputenc}
\usepackage{lmodern}
\usepackage{booktabs}
\usepackage{tabularx}
\usepackage{graphicx}
\usepackage{hyperref}
\usepackage{listings}
\usepackage{xcolor}
\usepackage{tikz}
\usetikzlibrary{positioning,arrows.meta}

% listings: Kotlin without shell-escape (no built-in Kotlin language)
\lstdefinelanguage{Kotlin}{
  morekeywords={
    abstract,actual,annotation,as,break,by,catch,class,companion,const,
    constructor,continue,crossinline,data,do,dynamic,else,enum,expect,
    external,false,field,file,final,finally,for,fun,get,if,import,in,
    infix,init,inline,inner,interface,internal,is,lateinit,noinline,null,
    object,open,operator,out,override,package,private,protected,public,
    reified,return,sealed,set,super,suspend,tailrec,this,throw,true,try,
    typealias,typeof,val,var,when,where,while
  },
  sensitive=true,
  morecomment=[l]{//},
  morecomment=[s]{/*}{*/},
  morestring=[b]",
  morestring=[b]',
  morestring=[s]{"""}{"""},
}

\lstdefinestyle{kotlinstyle}{
  language=Kotlin,
  basicstyle=\ttfamily\small,
  keywordstyle=\color{KeywordBlue}\bfseries,
  commentstyle=\color{CommentGray}\itshape,
  stringstyle=\color{StringGreen},
  numbers=left,
  numberstyle=\tiny\color{CommentGray},
  numbersep=6pt,
  frame=single,
  rulecolor=\color{LectureGray!40},
  backgroundcolor=\color{CodeBg},
  breaklines=true,
  showstringspaces=false,
  tabsize=2,
  captionpos=b,
  morekeywords={readln,readLine,println,print,listOf,mutableListOf,toInt,toDouble,toLong,toShort,toByte,toChar,toString,split,size},
}
\lstset{style=kotlinstyle}

\newcommand{\kotlincode}[1]{\texttt{\small #1}}
\newcommand{\trythis}[1]{%
  \begin{alertblock}{Try this}
    #1
  \end{alertblock}%
}
\newcommand{\commonmistake}[1]{%
  \begin{exampleblock}{Common mistake}
    #1
  \end{exampleblock}%
}
% Formative in-class checks: question slide, then a dedicated Answer slide
\newcommand{\quizpause}{%
  \par\vspace{0.5em}
  {\small\textit{Pause --- answer (clicker / raise hand / neighbour) before the next slide.}}%
}
\newcommand{\quizbox}[1]{%
  \setbeamercolor{block title}{bg=KeywordBlue,fg=white}%
  \setbeamercolor{block body}{bg=QuizBg,fg=black}%
  \begin{block}{Quiz}
    #1
  \end{block}%
  \setbeamercolor{block title}{bg=LectureNavy,fg=white}%
  \setbeamercolor{block body}{bg=CodeBg,fg=black}%
}
\newcommand{\answerbox}[1]{%
  \setbeamercolor{block title}{bg=AnswerGold,fg=white}%
  \setbeamercolor{block body}{bg=AnswerKeyBg,fg=black}%
  \begin{block}{Answer}
    #1
  \end{block}%
  \setbeamercolor{block title}{bg=LectureNavy,fg=white}%
  \setbeamercolor{block body}{bg=CodeBg,fg=black}%
}

% Empty author/date placeholders for the lecturer to fill
\author{}
\date{}
